% GENERATED FILE. Edit canonical lessons, then run npm run generate:books.
% canonical-source-hash: fnv1a64:475f5b7dcddde1b7
% canonical-lessons: PA-C246-lakkar, PA-C246-kari, PA-C246-chhappar, PA-C246-parnala, PA-C246-parchhatti

\chapter{Beams and Roofs}
\label{ch:pa-a2-beams-and-roofs}
\hlchaptermodality{246}

\begin{chapteropening}
\textbf{By the end of this chapter:} \emph{I can name what holds up a roof.}

Complete the last lesson of chapter 246: I can name what holds up a roof.
\end{chapteropening}

\section[\texorpdfstring{\pa{ਲੱਕੜ} (lakkaṛ)}{ਲੱਕੜ (lakkaṛ)}]{\pa{ਲੱਕੜ} (lakkaṛ) — wood}

\label{lesson:PA-C246-lakkar}



\begin{quote}
Before the new one: say the Punjabi for a pillar, then the Punjabi for a brick.
\end{quote}

\subsection*{You'll want to know: \pa{ਲੱਕੜ}}
\textbf{\pa{ਲੱਕੜ}} — \emph{lakkaṛ} — \textquotedblleft{}wood\textquotedblright{}.

\begin{quote}
\pa{ਇਹ} \pa{ਮੇਜ਼} \pa{ਲੱਕੜ} \pa{ਦਾ} \pa{ਹੈ।}
\end{quote}

\begin{quote}
\emph{eh mez lakkaṛ dā hai.}
\end{quote}

\begin{quote}
This table is made of wood.
\end{quote}

A feminine noun.

\begin{cousinweb}
\textbf{\pa{ਲੱਕੜ}} (\emph{lakkaṛ}) is inherited from Sanskrit \emph{lakuṭa}, \textquotedblleft{}club\textquotedblright{}; CDIAL lists Punjabi \emph{lakkaṛ}, \textquotedblleft{}wood\textquotedblright{} (CDIAL 10875). CDIAL marks \emph{lakkaṛ} masculine and lists a feminine Punjabi form, also \textquotedblleft{}wood\textquotedblright{}, beside it; here \pa{ਲੱਕੜ} is taught as feminine (\pa{ਗਿੱਲੀ} \pa{ਲੱਕੜ}).
\end{cousinweb}

\begin{grammarlens}[title={What you've built}]
One more word for this chapter.
\end{grammarlens}

\subsection*{Your turn}
\begin{itemize}
\raggedright
  \item \emph{Say it:} \emph{lakkaṛ}
  \item \emph{Say it:} \emph{lakkaṛ}, once more
  \item \emph{Say it:} \emph{eh mez lakkaṛ dā hai.} — This table is made of wood.
  \item \emph{Recall:} say the Punjabi for a pillar, then the Punjabi for a brick, then say \emph{lakkaṛ} again
\end{itemize}

\begin{tcolorbox}[breakable,colback=teal!4,colframe=teal!35!black,title={Before you move on}]
What does \pa{ਲੱਕੜ} mean? (\textquotedblleft{}Wood\textquotedblright{}.) Say it once more.
\end{tcolorbox}
\section[\texorpdfstring{\pa{ਕੜੀ} (kaṛī)}{ਕੜੀ (kaṛī)}]{\pa{ਕੜੀ} (kaṛī) — a roof beam; a rafter}

\label{lesson:PA-C246-kari}



\begin{quote}
Before the new one: say the Punjabi for a brick, then the Punjabi for wood.
\end{quote}

\subsection*{You'll want to know: \pa{ਕੜੀ}}
\textbf{\pa{ਕੜੀ}} — \emph{kaṛī} — \textquotedblleft{}a roof beam; a rafter\textquotedblright{}.

\begin{quote}
\pa{ਛੱਤ} \pa{ਦੀ} \pa{ਹਰ} \pa{ਕੜੀ} \pa{ਲੱਕੜ} \pa{ਦੀ} \pa{ਹੈ।}
\end{quote}

\begin{quote}
\emph{chhatt dī har kaṛī lakkaṛ dī hai.}
\end{quote}

\begin{quote}
Every beam of the roof is made of wood.
\end{quote}

A feminine noun.

\begin{cousinweb}
\textbf{\pa{ਕੜੀ}} (\emph{kaṛī}) is inherited from Sanskrit \emph{kaṭa}, \textquotedblleft{}thin piece of wood, plank\textquotedblright{}; CDIAL lists Punjabi \emph{kaṛī}, \textquotedblleft{}rafter\textquotedblright{} (CDIAL 2631).
\end{cousinweb}

\begin{grammarlens}[title={What you've built}]
One more word for this chapter.
\end{grammarlens}

\subsection*{Your turn}
\begin{itemize}
\raggedright
  \item \emph{Say it:} \emph{kaṛī}
  \item \emph{Say it:} \emph{kaṛī}, once more
  \item \emph{Say it:} \emph{chhatt dī har kaṛī lakkaṛ dī hai.} — Every beam of the roof is made of wood.
  \item \emph{Recall:} say the Punjabi for a brick, then the Punjabi for wood, then say \emph{kaṛī} again
  \item \emph{Write it:} copy \textbf{\pa{ਕੜੀ}} once, checking the whole word
\end{itemize}

\begin{tcolorbox}[breakable,colback=teal!4,colframe=teal!35!black,title={Before you move on}]
What does \pa{ਕੜੀ} mean? (\textquotedblleft{}A roof beam; a rafter\textquotedblright{}.) Say it once more.
\end{tcolorbox}
\section[\texorpdfstring{\pa{ਛੱਪਰ} (chhappar)}{ਛੱਪਰ (chhappar)}]{\pa{ਛੱਪਰ} (chhappar) — a thatched roof; a shed}

\label{lesson:PA-C246-chhappar}



\begin{quote}
Before the new one: say the Punjabi for wood, then the Punjabi for a rafter.
\end{quote}

\subsection*{You'll want to know: \pa{ਛੱਪਰ}}
\textbf{\pa{ਛੱਪਰ}} — \emph{chhappar} — \textquotedblleft{}a thatched roof; a shed\textquotedblright{}.

\begin{quote}
\pa{ਖੇਤ} \pa{ਵਿੱਚ} \pa{ਇੱਕ} \pa{ਛੋਟਾ} \pa{ਛੱਪਰ} \pa{ਹੈ।}
\end{quote}

\begin{quote}
\emph{khet vich ikk chhoṭā chhappar hai.}
\end{quote}

\begin{quote}
There is a small thatched shed in the field.
\end{quote}

A masculine noun.

\begin{cousinweb}
\textbf{\pa{ਛੱਪਰ}} (\emph{chhappar}) is inherited from Sanskrit \emph{chattvara}, \textquotedblleft{}house, bower\textquotedblright{}; CDIAL lists Punjabi \emph{chappar}, \textquotedblleft{}thatched roof\textquotedblright{} (CDIAL 4976).
\end{cousinweb}

\begin{grammarlens}[title={What you've built}]
One more word for this chapter.
\end{grammarlens}

\subsection*{Your turn}
\begin{itemize}
\raggedright
  \item \emph{Say it:} \emph{chhappar}
  \item \emph{Say it:} \emph{chhappar}, once more
  \item \emph{Say it:} \emph{khet vich ikk chhoṭā chhappar hai.} — There is a small thatched shed in the field.
  \item \emph{Recall:} say the Punjabi for wood, then the Punjabi for a rafter, then say \emph{chhappar} again
  \item \emph{Write it:} copy \textbf{\pa{ਛੱਪਰ}} once, checking the whole word
\end{itemize}

\begin{tcolorbox}[breakable,colback=teal!4,colframe=teal!35!black,title={Before you move on}]
What does \pa{ਛੱਪਰ} mean? (\textquotedblleft{}A thatched roof; a shed\textquotedblright{}.) Say it once more.
\end{tcolorbox}
\section[\texorpdfstring{\pa{ਪਰਨਾਲਾ} (parnālā)}{ਪਰਨਾਲਾ (parnālā)}]{\pa{ਪਰਨਾਲਾ} (parnālā) — a roof gutter; a rain spout}

\label{lesson:PA-C246-parnala}



\begin{quote}
Before the new one: say the Punjabi for a rafter, then the Punjabi for a thatched roof.
\end{quote}

\subsection*{You'll want to know: \pa{ਪਰਨਾਲਾ}}
\textbf{\pa{ਪਰਨਾਲਾ}} — \emph{parnālā} — \textquotedblleft{}a roof gutter; a rain spout\textquotedblright{}.

\begin{quote}
\pa{ਮੀਂਹ} \pa{ਵਿੱਚ} \pa{ਪਰਨਾਲੇ} \pa{ਤੋਂ} \pa{ਪਾਣੀ} \pa{ਆਉਂਦਾ} \pa{ਹੈ।}
\end{quote}

\begin{quote}
\emph{mīṁh vich parnāle ton pāṇī āundā hai.}
\end{quote}

\begin{quote}
When it rains, water comes down the gutter.
\end{quote}

A masculine noun.

\begin{cousinweb}
\textbf{\pa{ਪਰਨਾਲਾ}} (\emph{parnālā}) is inherited from Sanskrit \emph{pranāḍī}, \textquotedblleft{}outlet pipe, drain\textquotedblright{}; CDIAL lists Punjabi \emph{parnālā}, \textquotedblleft{}gutter, drain\textquotedblright{} (CDIAL 8673).
\end{cousinweb}

\begin{grammarlens}[title={What you've built}]
One more word for this chapter.
\end{grammarlens}

\subsection*{Your turn}
\begin{itemize}
\raggedright
  \item \emph{Say it:} \emph{parnālā}
  \item \emph{Say it:} \emph{parnālā}, once more
  \item \emph{Say it:} \emph{mīṁh vich parnāle ton pāṇī āundā hai.} — When it rains, water comes down the gutter.
  \item \emph{Recall:} say the Punjabi for a rafter, then the Punjabi for a thatched roof, then say \emph{parnālā} again
  \item \emph{Write it:} copy \textbf{\pa{ਪਰਨਾਲਾ}} once, checking the whole word
\end{itemize}

\begin{tcolorbox}[breakable,colback=teal!4,colframe=teal!35!black,title={Before you move on}]
What does \pa{ਪਰਨਾਲਾ} mean? (\textquotedblleft{}A roof gutter; a rain spout\textquotedblright{}.) Say it once more.
\end{tcolorbox}
\section[\texorpdfstring{\pa{ਪਰਛੱਤੀ} (parchhattī)}{ਪਰਛੱਤੀ (parchhattī)}]{\pa{ਪਰਛੱਤੀ} (parchhattī) — a high shelf; a shelf under the roof}

\label{lesson:PA-C246-parchhatti}



\begin{quote}
Before the new one: say the Punjabi for a thatched roof, then the Punjabi for a roof gutter.
\end{quote}

\subsection*{You'll want to know: \pa{ਪਰਛੱਤੀ}}
\textbf{\pa{ਪਰਛੱਤੀ}} — \emph{parchhattī} — \textquotedblleft{}a high shelf; a shelf under the roof\textquotedblright{}.

\begin{quote}
\pa{ਪੁਰਾਣੇ} \pa{ਭਾਂਡੇ} \pa{ਪਰਛੱਤੀ} \pa{ਉੱਤੇ} \pa{ਹਨ।}
\end{quote}

\begin{quote}
\emph{purāṇe bhāṇḍe parchhattī utte han.}
\end{quote}

\begin{quote}
The old pots are on the high shelf.
\end{quote}

A feminine noun.

\begin{cousinweb}
\textbf{\pa{ਪਰਛੱਤੀ}} (\emph{parchhattī}) is inherited from a reconstructed \emph{*praticchattika}, \textquotedblleft{}near the roof\textquotedblright{}; CDIAL lists Punjabi \emph{paṛchattī}, \textquotedblleft{}shelf under roof\textquotedblright{} (CDIAL 8556).
\end{cousinweb}

\begin{grammarlens}[title={What you've built}]
One more word for this chapter.
\end{grammarlens}

\subsection*{Your turn}
\begin{itemize}
\raggedright
  \item \emph{Say it:} \emph{parchhattī}
  \item \emph{Say it:} \emph{parchhattī}, once more
  \item \emph{Say it:} \emph{purāṇe bhāṇḍe parchhattī utte han.} — The old pots are on the high shelf.
  \item \emph{Recall:} say the Punjabi for a thatched roof, then the Punjabi for a roof gutter, then say \emph{parchhattī} again
  \item \emph{Write it:} copy \textbf{\pa{ਪਰਛੱਤੀ}} once, checking the whole word
\end{itemize}

\begin{tcolorbox}[breakable,colback=teal!4,colframe=teal!35!black,title={Before you move on}]
What does \pa{ਪਰਛੱਤੀ} mean? (\textquotedblleft{}A high shelf; a shelf under the roof\textquotedblright{}.) Say it once more.
\end{tcolorbox}
